\textbf{Claude 4.5 Sonnet} employed level-synchronous BFS (Figure~\ref{fig:bfs-level-sync}) across all languages, with varying degrees of success. The core strategy partitioned frontier vertices among workers, collected neighbors into private buffers, and merged results to form the next frontier. However, implementation details, particularly how merging preserved or violated discovery order, determined correctness and performance outcomes.

In Python and Go, the implementations followed GPT-5's approach but failed to achieve speedups. Python's \texttt{ProcessPoolExecutor} timed out on the dense benchmark due to inter-process communication overhead, while completing slowly on smaller test graphs. Go's goroutine-based version ran $3.3\times$
slower than sequential. Claude attributed this to inherent BFS limitations: level synchronization creates hard dependencies, small frontiers limit parallel work, and graph traversal is memory-bound. Notably, Claude explicitly identified these limitations rather than fabricating results.


Both implementations also introduced correctness violations through sorting. Python sorted the frontier by vertex ID at each level, producing numerically ordered output rather than preserving insertion order. Go avoided vertex-ID sorting by merging in chunk order, but normalized adjacency lists at startup with \texttt{g.Normalize()}, meaning output reflects sorted adjacency lists rather than original insertion order. Neither violation was caught by the respective test suites.

The C++ and Rust implementations attempted a two-phase approach: parallel neighbor collection followed by sequential merging. C++ used OpenMP with per-vertex buffers but ran 1.39$\times$ slower---the parallelized portion (copying neighbor lists) was too trivial to offset thread overhead. Rust achieved an 11.4$\times$ speedup on dense graphs using Rayon, but introduced \texttt{sort\_unstable()} at two points (the current level before appending to result, and the next frontier before the next iteration)---the same pattern appears in the submitted sequential baseline, so both produce numerically sorted output. This diverges from the reference insertion-order sequential BFS, causing 1 correctness failure on the unordered edge insertion test (24/25 tests pass).

The Java implementation was the only fully successful parallelization. Using a two-phase level-synchronous approach, workers collected neighbors into a fixed-position array (\texttt{perVertexCandidates[i]} for frontier vertex $i$), then a sequential pass processed buffers in order---preserving exact discovery order without sorting. With fallback thresholds (100 vertices for small graphs, 4 vertices for small frontiers), it achieved a 4.65$\times$ speedup on dense graphs and 2.1$\times$ on deep graphs---the only implementation to improve deep graph performance.

For C\#, Claude declined to parallelize entirely. After analyzing BFS constraints---level dependencies, discovery-order requirements, and visited-set synchronization---the model concluded that matching sequential output while parallelizing would effectively serialize most work. Rather than produce incorrect code or fabricate results, Claude returned the sequential algorithm with documented rationale, representing the most conservative response among all models tested.
